% Introduction (one sentence per line).
\section{Introduction}
\label{sec:intro}

Many extensions of the Standard Model predict particles that live long enough to travel a measurable distance before they decay~\cite{Dubois2008LLP,Vargas2024HNL}.
Their charged decay products start far from the beam line, so their tracks do not point back to the collision point~\cite{Kowalski2013Displaced}.
Searches for such signatures have become a central part of the physics programme of the LHC experiments.
Heavy neutral leptons and dark photons are two well-motivated examples~\cite{Vargas2024HNL,Chen2012DarkPhoton}.
In both cases the signal is a few tracks that appear in the middle of the detector.
Finding these tracks is the first step of every such search, and every later step inherits its losses.

Standard track reconstruction starts from seeds near the beam line and extends them outward with a Kalman filter~\cite{Hansen2014Kalman,Okafor2019Combinatorial}.
This strategy is efficient for prompt particles and fast enough for the trigger.
It loses particles whose first hit lies beyond the innermost layers, because no seed can be formed for them.
Dedicated large-radius tracking recovers part of these particles at a large cost in time and fake tracks~\cite{Rossi2018LRT}\footnote{Large-radius passes usually run only on the hits that the prompt pass left unused.}.
The cost grows with pile-up, which will increase several-fold at the high-luminosity LHC~\cite{Jensen2017Pileup,Okafor2013Upgrade}.
A reconstruction that treats prompt and displaced particles alike would remove this trade-off.

Early searches for long-lived particles relied on hand-written selections of displaced vertices.
Their reach was limited by the tracking, because a vertex can only be built from tracks that were found.
Later searches added dedicated passes of track finding with relaxed requirements on the impact parameter.
These passes raised the efficiency for displaced particles but also the number of fake tracks, which then had to be removed by tighter vertex selections.
The net gain depended strongly on the signal model, because the relaxed passes were tuned for particular lifetimes.

Graph neural networks offer such a reconstruction~\cite{Moreau2016GNNReview}.
They represent the hits of an event as the nodes of a graph and the possible track segments as its edges.
A network then learns which edges belong to the same particle~\cite{Fischer2011Interaction,Lindqvist2020MessagePassing}.
The Exa.TrkX pipeline showed that this approach reaches high efficiency on the TrackML detector~\cite{Yilmaz2023ExaTrkX,Tanaka2021TrackML}.
Because the graph does not assume that a particle comes from the beam line, displaced particles are not excluded by construction.
Recent studies extended these pipelines to large-radius tracks and found efficiencies close to those of prompt tracks~\cite{Sato2008EdgeClass,Zhang2013Metric}.

Learned tracking has developed quickly over the last decade.
The first studies classified hit pairs with simple networks and built tracks by hand.
Graph networks then made it possible to reason over many hits at once.
Metric learning replaced hand-written geometric rules in the construction of the graph.
End-to-end pipelines now take raw hits and return tracks with parameters.
Throughout this development, the training events were drawn from the process under study.

Learned trackers change this picture in two ways.
They can be trained to find displaced and prompt particles in one pass, without separate settings.
They can also be retrained cheaply when the detector or the run conditions change.
Both properties are attractive for a search programme that covers many models.
Both, however, rest on the assumption that the training events teach the network something general about tracks.

However, these studies trained the network on the signal they later evaluate.
A tracker trained on one model of new physics may favour that model and miss others.
In a search this is a bias: the selection would depend on the hypothesis being tested.
The machine-learning literature calls the question out-of-distribution generalization~\cite{Tanaka2015Generalization,Abe2020Domain}.
For tracking it has not been answered, because the training and test samples were always drawn from the same process.
We ask whether a tracker trained only on Standard Model events can find displaced particles as well as one trained on the signal.

This paper tests that assumption directly.
We do not propose a new architecture, and we do not tune the pipeline for displaced particles.
We take an established pipeline, keep every part of it fixed, and change only what it is trained on.
The comparison is therefore a measurement of what the training sample contributes.
A positive answer would mean that experiments can train one tracker on the abundant generic events they already simulate.

To answer it, we train two copies of the same pipeline.
One copy sees only top-quark pair events, which contain almost no displaced particles.
The other copy sees the displaced signal it will be tested on.
Both copies share the graph construction, the network, the track builder and every hyperparameter.
We then compare them on displaced particles generated after the design of the study was frozen.
A tracker that wins this comparison does not need the signal to be trained.

The answer is yes: the generic-trained tracker finds 2.49 points more displaced targets than the signal-trained one.
It also produces fewer fake tracks and fewer duplicates~\cite{Jensen2011Calibration}.
The result holds across particle species, displacement and pile-up.
It does not depend on the choice of training seed or on the bootstrap of the test events~\cite{Novak2023Bootstrap,Usman2011Seeds}.
The rest of this paper explains how we reached it.

Our contribution is a controlled measurement, not a new method.
It isolates the effect of the training sample on a learned tracker.
It reports that effect with an uncertainty that includes the randomness of the training.
And it checks that the effect survives changes of the signal model, the conditions and the analysis choices.

Section~\ref{sec:method} describes the detector, the samples, the graph and the network.
Section~\ref{sec:results} compares the two trackers and shows where the difference comes from.
Section~\ref{sec:summary} summarizes the result, and Appendix~\ref{app:stat} gives the details of the uncertainty.
The figures that carry the evidence are Fig.~\ref{fig:eff_disp} for the displacement and Fig.~\ref{fig:pileup} for pile-up.

Our study builds on a long line of work on fast tracking~\cite{Xu2010Hough,Esposito2015Cellular,Quinn2022Trigger}.
Graph-based pipelines have been ported to accelerators and now run within the time budget of a software trigger~\cite{Ito2021GPU,Petrov2009Inference}.
Their robustness to detector conditions has been studied for prompt particles~\cite{Hansen2008Robustness}.
Displaced particles have received less attention, because they are rare in the samples used for training.
This scarcity is exactly what makes generic training attractive: the network cannot learn the signal it never sees\footnote{Heavy-flavour hadrons travel a few millimetres before they decay, so the generic sample does contain some displaced tracks.}.

Filtering the candidate edges before classification is a common way to reduce the size of the graph~\cite{Garcia2018Filtering}.
We do not filter, because a filter trained on one sample could itself introduce the bias we want to measure.
The graph is therefore larger than in pipelines optimized for speed, and the edge classifier sees every candidate edge.
This choice makes the comparison cleaner at the cost of some time per event.

We simulate collisions at a centre-of-mass energy of 14~TeV with a generic all-silicon detector~\cite{Abe2009ACTS,Dubois2019Pythia}.
The detector response is simulated with a fast model that keeps the material and the magnetic field of the full geometry~\cite{Weber2014Geant,Kowalski2024Delphes}.
The study therefore measures the tracker, not the detector, and its numbers hold within this simulation only~\cite{Yilmaz2017Material,Rossi2012Resolution}.
We expect the comparison between the two trackers to carry over to a real detector, because both suffer the same effects.

The question matters beyond the LHC, where future colliders will face higher rates and new long-lived signatures~\cite{Chen2023FCC,Vargas2018Software}.
It also matters for model-independent searches, which need a reconstruction that does not encode a signal hypothesis~\cite{Moreau2010Anomaly}.
Statistical tools for such searches assume that the selection does not depend on the model being tested~\cite{Quinn2016Statistics}.
A generic-trained tracker satisfies this assumption by construction.

We keep the comparison fair in three ways.
First, both trackers use the same architecture and the same number of training events.
Second, every choice of the pipeline was fixed on validation events before the test events were examined.
Third, the test particles come from a signal sample that neither tracker was trained on.
With these precautions, a difference between the trackers can only come from what they were trained on.
